% Options for packages loaded elsewhere
\PassOptionsToPackage{unicode}{hyperref}
\PassOptionsToPackage{hyphens}{url}
\PassOptionsToPackage{dvipsnames,svgnames,x11names}{xcolor}
%
\documentclass[
  11pt,
  a4paper,
]{article}
\usepackage{amsmath,amssymb}
\usepackage{iftex}
\ifPDFTeX
  \usepackage[T1]{fontenc}
  \usepackage[utf8]{inputenc}
  \usepackage{textcomp} % provide euro and other symbols
\else % if luatex or xetex
  \usepackage{unicode-math} % this also loads fontspec
  \defaultfontfeatures{Scale=MatchLowercase}
  \defaultfontfeatures[\rmfamily]{Ligatures=TeX,Scale=1}
\fi
\usepackage{lmodern}
\ifPDFTeX\else
  % xetex/luatex font selection
\fi
% Use upquote if available, for straight quotes in verbatim environments
\IfFileExists{upquote.sty}{\usepackage{upquote}}{}
\IfFileExists{microtype.sty}{% use microtype if available
  \usepackage[]{microtype}
  \UseMicrotypeSet[protrusion]{basicmath} % disable protrusion for tt fonts
}{}
\makeatletter
\@ifundefined{KOMAClassName}{% if non-KOMA class
  \IfFileExists{parskip.sty}{%
    \usepackage{parskip}
  }{% else
    \setlength{\parindent}{0pt}
    \setlength{\parskip}{6pt plus 2pt minus 1pt}}
}{% if KOMA class
  \KOMAoptions{parskip=half}}
\makeatother
\usepackage{xcolor}
\usepackage[margin=25mm]{geometry}
\usepackage{color}
\usepackage{fancyvrb}
\newcommand{\VerbBar}{|}
\newcommand{\VERB}{\Verb[commandchars=\\\{\}]}
\DefineVerbatimEnvironment{Highlighting}{Verbatim}{commandchars=\\\{\}}
% Add ',fontsize=\small' for more characters per line
\newenvironment{Shaded}{}{}
\newcommand{\AlertTok}[1]{\textcolor[rgb]{1.00,0.00,0.00}{\textbf{#1}}}
\newcommand{\AnnotationTok}[1]{\textcolor[rgb]{0.38,0.63,0.69}{\textbf{\textit{#1}}}}
\newcommand{\AttributeTok}[1]{\textcolor[rgb]{0.49,0.56,0.16}{#1}}
\newcommand{\BaseNTok}[1]{\textcolor[rgb]{0.25,0.63,0.44}{#1}}
\newcommand{\BuiltInTok}[1]{\textcolor[rgb]{0.00,0.50,0.00}{#1}}
\newcommand{\CharTok}[1]{\textcolor[rgb]{0.25,0.44,0.63}{#1}}
\newcommand{\CommentTok}[1]{\textcolor[rgb]{0.38,0.63,0.69}{\textit{#1}}}
\newcommand{\CommentVarTok}[1]{\textcolor[rgb]{0.38,0.63,0.69}{\textbf{\textit{#1}}}}
\newcommand{\ConstantTok}[1]{\textcolor[rgb]{0.53,0.00,0.00}{#1}}
\newcommand{\ControlFlowTok}[1]{\textcolor[rgb]{0.00,0.44,0.13}{\textbf{#1}}}
\newcommand{\DataTypeTok}[1]{\textcolor[rgb]{0.56,0.13,0.00}{#1}}
\newcommand{\DecValTok}[1]{\textcolor[rgb]{0.25,0.63,0.44}{#1}}
\newcommand{\DocumentationTok}[1]{\textcolor[rgb]{0.73,0.13,0.13}{\textit{#1}}}
\newcommand{\ErrorTok}[1]{\textcolor[rgb]{1.00,0.00,0.00}{\textbf{#1}}}
\newcommand{\ExtensionTok}[1]{#1}
\newcommand{\FloatTok}[1]{\textcolor[rgb]{0.25,0.63,0.44}{#1}}
\newcommand{\FunctionTok}[1]{\textcolor[rgb]{0.02,0.16,0.49}{#1}}
\newcommand{\ImportTok}[1]{\textcolor[rgb]{0.00,0.50,0.00}{\textbf{#1}}}
\newcommand{\InformationTok}[1]{\textcolor[rgb]{0.38,0.63,0.69}{\textbf{\textit{#1}}}}
\newcommand{\KeywordTok}[1]{\textcolor[rgb]{0.00,0.44,0.13}{\textbf{#1}}}
\newcommand{\NormalTok}[1]{#1}
\newcommand{\OperatorTok}[1]{\textcolor[rgb]{0.40,0.40,0.40}{#1}}
\newcommand{\OtherTok}[1]{\textcolor[rgb]{0.00,0.44,0.13}{#1}}
\newcommand{\PreprocessorTok}[1]{\textcolor[rgb]{0.74,0.48,0.00}{#1}}
\newcommand{\RegionMarkerTok}[1]{#1}
\newcommand{\SpecialCharTok}[1]{\textcolor[rgb]{0.25,0.44,0.63}{#1}}
\newcommand{\SpecialStringTok}[1]{\textcolor[rgb]{0.73,0.40,0.53}{#1}}
\newcommand{\StringTok}[1]{\textcolor[rgb]{0.25,0.44,0.63}{#1}}
\newcommand{\VariableTok}[1]{\textcolor[rgb]{0.10,0.09,0.49}{#1}}
\newcommand{\VerbatimStringTok}[1]{\textcolor[rgb]{0.25,0.44,0.63}{#1}}
\newcommand{\WarningTok}[1]{\textcolor[rgb]{0.38,0.63,0.69}{\textbf{\textit{#1}}}}
\usepackage{longtable,booktabs,array}
\usepackage{calc} % for calculating minipage widths
% Correct order of tables after \paragraph or \subparagraph
\usepackage{etoolbox}
\makeatletter
\patchcmd\longtable{\par}{\if@noskipsec\mbox{}\fi\par}{}{}
\makeatother
% Allow footnotes in longtable head/foot
\IfFileExists{footnotehyper.sty}{\usepackage{footnotehyper}}{\usepackage{footnote}}
\makesavenoteenv{longtable}
\setlength{\emergencystretch}{3em} % prevent overfull lines
\providecommand{\tightlist}{%
  \setlength{\itemsep}{0pt}\setlength{\parskip}{0pt}}
\setcounter{secnumdepth}{-\maxdimen} % remove section numbering
\ifLuaTeX
\usepackage[bidi=basic]{babel}
\else
\usepackage[bidi=default]{babel}
\fi
\babelprovide[main,import]{english}
% get rid of language-specific shorthands (see #6817):
\let\LanguageShortHands\languageshorthands
\def\languageshorthands#1{}
\usepackage{mathpazo}
\usepackage{mathrsfs}
\usepackage{microtype}
\usepackage{fancyhdr}
\usepackage{needspace}
\usepackage{titlesec}
\usepackage{fvextra}
\DefineVerbatimEnvironment{verbatim}{Verbatim}{breaklines,fontsize=\small}
\pagestyle{fancy}
\fancyhf{}
\fancyhead[L]{\small\textsc{Depth-Generated Geometry}}
\fancyhead[R]{\small Research Snapshot 2026-09-07}
\fancyfoot[C]{\thepage}
\setlength{\headheight}{14pt}
\renewcommand{\headrulewidth}{0.3pt}
\setlength{\parskip}{5pt}
\setlength{\parindent}{0pt}
\setlength{\emergencystretch}{3em}
\titleformat{\section}{\large\bfseries}{\thesection}{0.6em}{}
\titlespacing*{\section}{0pt}{18pt}{8pt}
\titleformat{\subsection}{\normalsize\bfseries}{\thesubsection}{0.6em}{}
\titlespacing*{\subsection}{0pt}{12pt}{6pt}
\let\originalsection\section
\renewcommand{\section}{\needspace{7\baselineskip}\originalsection}
\allowdisplaybreaks[1]
\AtBeginDocument{\hypersetup{linkcolor=black,urlcolor=blue}}
\ifLuaTeX
  \usepackage{selnolig}  % disable illegal ligatures
\fi
\IfFileExists{bookmark.sty}{\usepackage{bookmark}}{\usepackage{hyperref}}
\IfFileExists{xurl.sty}{\usepackage{xurl}}{} % add URL line breaks if available
\urlstyle{same}
\hypersetup{
  pdftitle={Ordered Quaternionic Response Quotients},
  pdfauthor={Residual Chart Lab},
  pdflang={en},
  colorlinks=true,
  linkcolor={Maroon},
  filecolor={Maroon},
  citecolor={Blue},
  urlcolor={Blue},
  pdfcreator={LaTeX via pandoc}}

\title{Ordered Quaternionic Response Quotients}
\usepackage{etoolbox}
\makeatletter
\providecommand{\subtitle}[1]{% add subtitle to \maketitle
  \apptocmd{\@title}{\par {\large #1 \par}}{}{}
}
\makeatother
\subtitle{Flat Cores, Local Residuals, and a Three-Spectator Descent
Obstruction}
\author{Residual Chart Lab}
\date{Depth-Generated Geometry - Research Snapshot 2026-09-07}

\begin{document}
\maketitle

\hypertarget{abstract}{%
\subsection{Abstract}\label{abstract}}

We study finite cokernels formed from four ordered quaternion-valued
face responses and their six pairwise common shadows. The incidence of
labelled edge images distinguishes quotients with the same abstract
dimension. The one-spectator family admits two explicit coordinate
charts with a central cross-product transition. At two internal
spectators, six reduced spacing words carry a flat 144-dimensional
stationary-edge core connection, while the word 212 has an additional
quaternionic residual invisible to its incident hinges. We then
construct the three-spectator quotient at 222. It has dimension 444, an
intrinsic 432-dimensional outer-edge subspace, and a 12-dimensional
residual quotient. A specified central-slot right decoder induces an
isomorphism from the 212 residual tensored with the vector
representation onto this new residual quotient. The same suspension
fails to descend to the complete quotients: its obstruction image is
exactly the 432-dimensional core. All finite rank claims are over the
rationals and are accompanied by exact arithmetic certificates. The
comparison operator is part of the theorem; no choice-independent full
transport or nonzero core curvature is asserted.

\textbf{Snapshot scope.} This research snapshot consolidates the results
for mathematical review. The immutable Free Numbers Core v1.0.0 artifact
is a cited foundation, not a version replaced by this snapshot. The
manuscript reorganizes established source notes and the Note 26 result;
it does not claim a complete three-spectator atlas or a physical
interpretation. This is revision research-2026-09-07-audit1; the
companion audit report records the corrections and the verification
boundary.

\hypertarget{overview}{%
\subsection{Overview}\label{overview}}

We first describe the finite atlas through two internal spectators, then
construct the residual comparison for the three-spectator word 222.

Four ordered face responses meet along six labelled edge responses. At
the minimal tetrahedral level their compatibility quotient is

\[
\mathbb H\otimes\mathbb H
=
W_{12}\mathbin{\overset\perp\oplus}K_4,
\qquad
\dim W_{12}=12,
\qquad
K_4\cong\mathbb H.
\tag{0.1}
\]

Insert one spectator among the five ordered gap positions. Each of the
five resulting local quotients still has the same abstract
48-dimensional type

\[
T_6=(\mathbb H\otimes\mathbb H)\otimes V,
\qquad V=\operatorname{Im}\mathbb H.
\tag{0.2}
\]

Nevertheless, no single direct seed coordinate covers all five
placements. Two chart families are necessary:

\[
U_L=\{1,2,3\},
\qquad
U_R=\{1,3,4,5\}.
\tag{0.3}
\]

Their exterior overlap is trivial, while their central overlap is not:

\[
G_1=I,
\qquad
G_3=\operatorname{id}_{\mathbb H}\otimes\theta.
\tag{0.4}
\]

For \(a,w\in V\),

\[
\boxed{
\begin{aligned}
\theta(1\otimes w)
&=1\otimes w+
\sum_{\rho=1}^{3}e_\rho\otimes(e_\rho\times w),\\
\theta(a\otimes w)&=-w\otimes a,
\end{aligned}}
\qquad
a\times w=\frac12(aw-wa).
\tag{0.5}
\]

Thus the abstract local module does not record the whole object. The
incidence of its labelled edge images and the transition between valid
local coordinates retain where the spectator was inserted.

The shortest accurate summary is therefore

\[
\boxed{
\text{same local quotient type}
\;\not\Rightarrow\;
\text{same placement coordinate},
}
\tag{0.6}
\]

and the first nontrivial coordinate change is an explicit cross-product
shear.

\begin{center}\rule{0.5\linewidth}{0.5pt}\end{center}

\hypertarget{minimal-setup}{%
\subsection{1. Minimal setup}\label{minimal-setup}}

Let

\[
\mathbb H=\operatorname{span}_{\mathbb R}\{1,i,j,k\},
\qquad
V=\operatorname{Im}\mathbb H.
\tag{1.1}
\]

For response depth \(m\), put

\[
\mathcal R_m=\operatorname{Hom}(V^{\otimes m},\mathbb H).
\tag{1.2}
\]

Choose four ordered gaps

\[
Q=\{q_1<q_2<q_3<q_4\}.
\tag{1.3}
\]

They are treated as the four vertices of an ordered tetrahedron. The six
local edges are labelled

\[
12,13,14,23,24,34.
\tag{1.4}
\]

Four face responses restrict to the six pairwise common shadows. Their
signed disagreement defines the matching map

\[
\partial_{n,Q}:
\mathcal R_{n-2}^{\oplus4}
\longrightarrow
\mathcal R_{n-3}^{\oplus6},
\tag{1.5}
\]

\[
(\partial_{n,Q}F)_{ab}
=
\rho_{q_a}^{q_aq_b}F_a
-
\rho_{q_b}^{q_aq_b}F_b.
\tag{1.6}
\]

The local relation object is the cokernel

\[
Y_{n,Q}:=\operatorname{coker}\partial_{n,Q}.
\tag{1.7}
\]

For each labelled edge, retain its image in the quotient:

\[
E_{ab}^{(n,Q)}
:=
\operatorname{im}
\left(
\mathcal R_{n-3}^{(ab)}\longrightarrow Y_{n,Q}
\right).
\tag{1.8}
\]

The restrictions in (1.6) have the following concrete construction. For
\(T=v_1\otimes\cdots\otimes v_n\) and a set of omitted gaps
\(S\subset\{1,\ldots,n-1\}\), set \[
 A_{n,S}(T)(x)=v_n b_{n-1}v_{n-1}\cdots b_1v_1,
 \qquad
 b_g=\begin{cases}1,&g\in S,\\x_g,&g\notin S.\end{cases}
\] Every terminal face \(A_{n,\{r\}}\) is onto. One can see this by
replacing \(v_{r+1}v_r\) by an arbitrary quaternion \(h\), then
successively using the invertible one-variable left and right encoders
on the remaining probes. The right encoder is \[
 \Theta_R(h\otimes w)(z)=hzw,
 \qquad
 J_m(h;v_1,\ldots,v_m)(x_1,\ldots,x_m)
 =h x_m v_m\cdots x_1v_1.
\] The left encoder is \(\Theta_L(w\otimes h)(z)=wzh\). Their matrices
in the quaternion basis are invertible. Thus the common-shadow
restriction, when factored through a face, is unique: \[
 \rho_r^{rs}A_{n,\{r\}}=A_{n,\{r,s\}}.
\] Its existence is the local decoder-collapse identity of Note 12. For
all finite supports reported here it is also checked entry by entry
against the literal product above. This fixes the matching map without
choosing a basis of its cokernel. It does not require the global
all-length filling theorem.

We work over the rational quaternion algebra with
\(i^2=j^2=k^2=ijk=-1\), and extend scalars to \(\mathbb R\) for the
usual diagonal \(SO(3)\)-action. All reported rational ranks therefore
also hold over \(\mathbb R\). The word \emph{atlas} refers here to a
discrete family of quotient coordinates and their labelled comparison
maps; no ambient manifold is assumed. Unless actual gap labels are
explicitly stated, edge labels \(ab\) denote the \(a\)-th and \(b\)-th
support vertices. Local \(14\), for example, is the long edge even when
its actual gap labels are \((1,6)\).

The object relevant to placement is not only \(Y_{n,Q}\), but the
decorated quotient

\[
\mathscr Y_{n,Q}
=
\left(Y_{n,Q};E_{12},E_{13},E_{14},E_{23},E_{24},E_{34}\right).
\tag{1.9}
\]

This distinction is the starting point of the atlas.

\begin{center}\rule{0.5\linewidth}{0.5pt}\end{center}

\hypertarget{the-seed-tetrahedron-at-n5}{%
\subsection{\texorpdfstring{2. The seed tetrahedron at
\(n=5\)}{2. The seed tetrahedron at n=5}}\label{the-seed-tetrahedron-at-n5}}

At \(n=5\), there is one four-gap support. Note 14 replaces an abstract
quotient basis by a closed, \(SO(3)\)-equivariant operator

\[
\omega_5:
\mathcal R_2^{\oplus6}
\longrightarrow
\mathbb H\otimes\mathbb H
\tag{2.1}
\]

satisfying

\[
\boxed{
\omega_5\partial_5=0,
\qquad
\operatorname{rank}\omega_5=16,
\qquad
\ker\omega_5=\operatorname{im}\partial_5.
}
\tag{2.2}
\]

Hence

\[
Y_{5,G_5}\cong\mathbb H\otimes\mathbb H.
\tag{2.3}
\]

Put \(e_0=1,e_1=i,e_2=j,e_3=k\). Define five maps

\[
P,A_0,B_0,C_0,D_0:
\mathbb H\otimes V\otimes V
\longrightarrow\mathbb H\otimes\mathbb H
\]

by

\[
\begin{aligned}
P(h,u,v)
&:=h\otimes uv,\\[1mm]
A_0(h,u,v)
&:=\sum_{a=0}^{3}he_a\otimes uv e_a,\\[1mm]
B_0(h,u,v)
&:=\sum_{a=0}^{3}he_a\otimes u e_a v,\\[1mm]
C_0(h,u,v)
&:=\sum_{a=0}^{3}e_a\otimes uvh e_a,\\[1mm]
D_0(h,u,v)
&:=\sum_{a=0}^{3}e_a h e_a v\otimes u.
\end{aligned}
\]

Associativity fixes all parentheses. Each map is \(SO(3)\)-equivariant:
quaternion multiplication is equivariant for conjugation, and every
summed pair of basis elements comes from the invariant tensor
\(\sum_{a=0}^3e_a\otimes e_a\).

The six ordered blocks are \[
\begin{aligned}
\Lambda_{12}&=P-\tfrac12B_0+\tfrac12C_0-D_0,\\
\Lambda_{13}&=P-\tfrac14A_0-\tfrac12B_0+\tfrac12C_0-D_0,\\
\Lambda_{14}&=-\tfrac14A_0-\tfrac12B_0,\\
\Lambda_{23}&=-\tfrac14A_0-\tfrac12C_0,\\
\Lambda_{24}&=P-\tfrac14A_0,\qquad \Lambda_{34}=P.
\end{aligned}
\] Thus \(\omega_5(F)=\sum_{a<b}\Lambda_{ab}J_2^{-1}F_{ab}\).
Substitution gives four signed vertex identities \[
\sum_{b>a}\Lambda_{ab}J_2^{-1}\rho_{q_a}^{q_aq_b}
=\sum_{b<a}\Lambda_{ba}J_2^{-1}\rho_{q_a}^{q_bq_a},
\] which prove \(\omega_5\partial_5=0\). The block \(\Lambda_{34}=P\) is
onto, and the matching matrix has rank 200 in its 216-dimensional
target. These finite identities and ranks are verified over exact
rational arithmetic by the seed certificate. They prove (2.2).

Write \(\Lambda_{ab}\) for the six edge blocks of \(\omega_5\). Their
ranks are

\[
\begin{array}{c|rrrrrr}
ab&12&13&14&23&24&34\\ \hline
\operatorname{rank}\Lambda_{ab}&16&12&16&4&12&16.
\end{array}
\tag{2.4}
\]

The two rank-12 images coincide, while the central rank-4 image supplies
a complement:

\[
W_{12}:=\operatorname{im}\Lambda_{13}
=\operatorname{im}\Lambda_{24},
\qquad
K_4:=\operatorname{im}\Lambda_{23},
\tag{2.5}
\]

\[
\mathbb H\otimes\mathbb H=W_{12}\oplus K_4.
\tag{2.6}
\]

Note 18 makes this split canonical. Define

\[
\iota(q)=\sum_{\alpha=0}^{3}e_\alpha\otimes qe_\alpha,
\qquad
\nu=\frac14\iota^*,
\tag{2.7}
\]

where \(e_0=1,e_1=i,e_2=j,e_3=k\). Then

\[
\boxed{
K_4=\iota(\mathbb H),
\qquad
W_{12}=\ker\nu=\iota(\mathbb H)^\perp.
}
\tag{2.8}
\]

Thus the seed contains a canonical orthogonal \(12+4\) channel split,
not merely a dimension decomposition.

The adjoint uses the Euclidean quaternion inner product and its tensor
product, with \(1,i,j,k\) orthonormal. Explicitly,
\(\iota^*(x\otimes y)=y\bar x\) and \(\iota^*\iota=4I\).

The quaternion coordinate of the six edge blocks has the exact incidence
pattern

\[
\begin{array}{c|rrrrrr}
ab&12&13&14&23&24&34\\ \hline
\operatorname{rank}(\nu\Lambda_{ab})&4&0&4&4&0&4\\
\text{outer sign}&+&0&-&\text{central}&0&+
\end{array}
\tag{2.9}
\]

This labelled pattern is the seed datum to be compared after a spectator
is inserted.

\begin{center}\rule{0.5\linewidth}{0.5pt}\end{center}

\hypertarget{one-spectator-and-five-placements}{%
\subsection{3. One spectator and five
placements}\label{one-spectator-and-five-placements}}

Let

\[
B_5=\{1,2,3,4,5\},
\qquad
Q_s=B_5\setminus\{s\}.
\tag{3.1}
\]

The omitted gap \(s\) is the spectator position. For every \(s\),

\[
\dim Y_{6,Q_s}=48,
\qquad
Y_{6,Q_s}\cong T_6
:=(\mathbb H\otimes\mathbb H)\otimes V.
\tag{3.2}
\]

The equality of abstract type does not produce one common seed
coordinate. Use the three full-rank seed edges as possible anchors:

\[
L=12,
\qquad
M=14,
\qquad
R=34.
\tag{3.3}
\]

Exact rational classification gives

\[
\begin{array}{c|c|ccc}
s&Q_s&L&M&R\\ \hline
1&(2,3,4,5)&\checkmark&\checkmark&\checkmark\\
2&(1,3,4,5)&\checkmark&\checkmark&-\\
3&(1,2,4,5)&\checkmark&-&\checkmark\\
4&(1,2,3,5)&-&-&\checkmark\\
5&(1,2,3,4)&-&-&\checkmark.
\end{array}
\tag{3.4}
\]

Every marked entry determines a unique surjective quotient coordinate

\[
\Omega_s^e:
(\mathbb H\otimes V^{\otimes3})^{\oplus6}
\longrightarrow T_6
\tag{3.5}
\]

whose kernel is exactly the image of the local matching map:

\[
\boxed{
\ker\Omega_s^e=\operatorname{im}\widehat\partial_{6,Q_s}.
}
\tag{3.6}
\]

The unmarked entries are impossible over \(\mathbb Q\), not merely
absent from one computation.

For precision, the direct normalization on the anchor edge \(e\) is \[
 (\Omega_s^e)_e(h;v_1,v_2,v_3)
 =\Lambda_e(h;v_{a_1},v_{a_2})\otimes v_t.
\] Here the three remaining probe slots are in increasing gap order;
\(a_1<a_2\) are the two slots belonging to the support and \(t\) is the
spectator slot. Compatibility uniquely completes each marked anchor.

The middle chart adds no new family. The two essential domains are

\[
\boxed{
U_L=\{1,2,3\},
\qquad
U_R=\{1,3,4,5\},
\qquad
U_L\cup U_R=B_5.
}
\tag{3.7}
\]

This is the first exact placement atlas in the program.

\begin{center}\rule{0.5\linewidth}{0.5pt}\end{center}

\hypertarget{the-overlap-transition}{%
\subsection{4. The overlap transition}\label{the-overlap-transition}}

The two charts overlap at one exterior placement and one central
placement:

\[
U_L\cap U_R=\{1,3\}.
\tag{4.1}
\]

At \(s=1\),

\[
\boxed{\Omega_1^L=\Omega_1^R.}
\tag{4.2}
\]

At \(s=3\),

\[
\boxed{
\Omega_3^L=G\Omega_3^R,
\qquad
G=\operatorname{id}_{\mathbb H}\otimes\theta.
}
\tag{4.3}
\]

Put

\[
\delta(w)=
\sum_{\rho=1}^{3}e_\rho\otimes(e_\rho\times w).
\tag{4.4}
\]

Then the transition is the closed map

\[
\boxed{
\begin{aligned}
\theta(1\otimes w)&=1\otimes w+\delta(w),\\
\theta(a\otimes w)&=-w\otimes a
\qquad(a,w\in V).
\end{aligned}}
\tag{4.5}
\]

The first line contains the essential alternating correction. A
spectator outside the active region gives no correction; a spectator
between the two active pairs forces a cross-product term.

The internal decomposition

\[
\mathbb H\otimes V
=
(1\otimes V)\oplus\delta(V)\oplus\operatorname{Sym}^2V
\tag{4.6}
\]

puts \(\theta\) into the form

\[
\theta=R+N,
\qquad
R^2=I,
\qquad
N^2=0,
\qquad
\operatorname{rank}N=3.
\tag{4.7}
\]

Its minimal polynomial is

\[
\boxed{m_\theta(t)=(t-1)^2(t+1).}
\tag{4.8}
\]

Therefore the transition is not a permutation of tensor slots. It is an
involutive reflection together with a three-dimensional nilpotent shear.
Over characteristic zero it has infinite order.

In (4.7), \(R\) is the identity on \(1\otimes V\) and minus the tensor
flip on \(V\otimes V\); \(N(1\otimes w)=\delta(w)\) and
\(N(V\otimes V)=0\). In particular \(RN=NR\).

This is a statement about a coordinate transition. No repeated power in
(4.8) is yet interpreted as a geometric loop.

\begin{center}\rule{0.5\linewidth}{0.5pt}\end{center}

\hypertarget{what-the-placement-remembers}{%
\subsection{5. What the placement
remembers}\label{what-the-placement-remembers}}

At the central support

\[
Q_3=(1,2,4,5),
\tag{5.1}
\]

the long edge \(14=(1,5)\) has image dimension 44 rather than 48. Its
four-dimensional quotient is detected by the cap residual of Note 17.

In the right chart define

\[
\beta_R((x\otimes y)\otimes w)=xw\bar y.
\tag{5.2}
\]

On actual responses define \[
 \lambda_3^L(F)=\sum_{a,b=1}^3 e_bF(e_a,e_a,e_b),\qquad
 \lambda_3^R(F)=\sum_{a,b=1}^3 F(e_b,e_a,e_a)e_b.
\] The chart in (3.5) uses decoded edge tensors, so write
\(\widehat\lambda_3^\bullet=\lambda_3^\bullet J_3\) and
\(\widehat\chi_6=\chi_6(\bigoplus_e J_3)\). This makes the domains in
the following identity explicit.

Then

\[
\boxed{
\beta_R\Omega_3^R
=
\widehat\chi_6
=
(-\widehat\lambda_3^L,+\widehat\lambda_3^L,0,0,
-\widehat\lambda_3^R,+\widehat\lambda_3^R).
}
\tag{5.3}
\]

Identify \(Y_{6,Q_3}\) with \(T_6\) by the induced right chart, so that
the edge images below are subspaces of \(T_6\). Then

\[
\boxed{
E_{14}=\ker\beta_R,
\qquad
T_6/E_{14}\cong\mathbb H.
}
\tag{5.4}
\]

Tensor the seed split (2.8) with \(V\):

\[
T_6
=
(W_{12}\otimes V)
\mathbin{\overset\perp\oplus}
(K_4\otimes V).
\tag{5.5}
\]

The direct extension of the seed quaternion does not survive the
long-edge quotient:

\[
\boxed{
\beta_R(K_4\otimes V)=0,
\qquad
K_4\otimes V\subset E_{14}.
}
\tag{5.6}
\]

The surviving quaternion is instead a quotient of \(W_{12}\otimes V\).
Thus the spectator does not carry the seed channel forward by a
placement-blind tensor product. It changes which channel represents the
residual.

The transition \(G\) explains this transfer. It mixes the two summands
in (5.5), and in particular

\[
G(K_4\otimes V)\cap(K_4\otimes V)=0.
\tag{5.7}
\]

The seed quaternion has not been destroyed. Its coordinate realization
has changed.

\begin{center}\rule{0.5\linewidth}{0.5pt}\end{center}

\hypertarget{the-residual-is-independent-of-the-chart}{%
\subsection{6. The residual is independent of the
chart}\label{the-residual-is-independent-of-the-chart}}

Transporting \(\beta_R\) across the central transition gives

\[
\beta_L:=\beta_RG^{-1},
\tag{6.1}
\]

with the closed formula

\[
\beta_L((x\otimes y)\otimes w)=x(2y+\bar y)w.
\tag{6.2}
\]

The same observable is recovered from either chart:

\[
\boxed{
\widehat\chi_6
=
\beta_R\Omega_3^R
=
\beta_L\Omega_3^L.
}
\tag{6.3}
\]

This separates two roles:

\begin{itemize}
\tightlist
\item
  \(G\) records how the local coordinate changes with placement;
\item
  \(\chi_6\) is the quaternionic residual unchanged by that coordinate
  change.
\end{itemize}

The coefficient \(2y+\bar y\) is not new bookkeeping. The same
scalar-vector weighting already occurs in the canonical central
coordinate of the \(n=5\) seed. The atlas transition forces it to
reappear.

\begin{center}\rule{0.5\linewidth}{0.5pt}\end{center}

\hypertarget{the-exceptional-two-spectator-quotient-at-n7}{%
\subsection{\texorpdfstring{7. The exceptional two-spectator quotient at
\(n=7\)}{7. The exceptional two-spectator quotient at n=7}}\label{the-exceptional-two-spectator-quotient-at-n7}}

The \(n=6\) atlas is complete for one spectator. Notes 16 and 21
identify and then fully coordinate the exceptional two-spectator
continuation.

The first ordered support with parity separation

\[
\mathrm O\,\mathrm O\mid\mathrm E\,\mathrm E
\tag{7.1}
\]

is uniquely

\[
\boxed{(1,3\mid4,6),}
\tag{7.2}
\]

whose spacing is \((2,1,2)\). On the four cross edges define

\[
\kappa_{212}
(F_{13},F_{14},F_{23},F_{24})
=
\varepsilon_4(F_{13})
-\varepsilon_4(F_{14})
-\varepsilon_4(F_{23})
+\varepsilon_4(F_{24}),
\tag{7.3}
\]

where

\[
\varepsilon_4(F)
=
\sum_{a,b=1}^{3}F(e_a,e_a,e_b,e_b).
\tag{7.4}
\]

Then, over \(\mathbb Q\),

\[
\boxed{
\mathcal R_5^{\oplus4}
\xrightarrow{\partial_\square}
\mathcal R_4^{\oplus4}
\xrightarrow{\kappa_{212}}
\mathbb H
\longrightarrow0,
\qquad
\ker\kappa_{212}=\operatorname{im}\partial_\square.
}
\tag{7.5}
\]

Thus the exceptional four-dimensional residual at \(n=7\) is the
complete cokernel of an alternating quaternionic cross square. Its
position is forced by order and parity, not guessed from the dimension
anomaly.

Note 21 transports a full-rank edge anchor from the \(n=6\) atlas and
proves that it has a unique compatible completion

\[
C_{212}:E_{(1,3,4,6)}\longrightarrow
T_7:=(\mathbb H\otimes\mathbb H)\otimes V^{\otimes2}.
\tag{7.6}
\]

The core coordinate and the square residual exactly exhaust the
quotient:

\[
\boxed{
Y_{7,(1,3,4,6)}
\cong T_7\oplus\mathbb H,
\qquad
\dim Y_{7,(1,3,4,6)}=144+4=148.
}
\tag{7.7}
\]

The six labelled edge ranks split as

\[
\begin{array}{c|rrrrrr}
&12&13&14&23&24&34\\ \hline
\text{core}&144&108&144&36&108&144\\
\text{full exceptional}&144&112&148&40&112&144.
\end{array}
\tag{7.8}
\]

Thus the extra quaternion is transverse to the transported
144-dimensional spectator core and is seen exactly on the four cross
edges.

\begin{center}\rule{0.5\linewidth}{0.5pt}\end{center}

\hypertarget{the-complete-reduced-two-spectator-atlas}{%
\subsection{8. The complete reduced two-spectator
atlas}\label{the-complete-reduced-two-spectator-atlas}}

Exterior suspension strips off every spectator lying strictly outside
the support interval. At \(n=7\), the irreducible two-spectator problem
therefore consists of the six positive spacing words

\[
113,\quad122,\quad131,\quad212,\quad221,\quad311.
\tag{8.1}
\]

For each of the five generic words \(\lambda\ne212\), Note 23 transports
one full-rank \(n=6\) edge anchor and proves that local compatibility
completes it uniquely to an exact quotient coordinate

\[
C_\lambda:E_{Q_\lambda}\longrightarrow T_7,
\qquad
\ker C_\lambda=\operatorname{im}\widehat\partial_{7,Q_\lambda}.
\tag{8.2}
\]

Consequently the full characteristic-zero dimension row is

\[
\boxed{
\begin{array}{c|rrrrrr}
\lambda&113&122&131&212&221&311\\ \hline
\dim Y_{7,Q_\lambda}&144&144&144&148&144&144.
\end{array}}
\tag{8.3}
\]

The generic words fall into two edge-incidence classes. The
exterior-type words \(113,311\) have ranks

\[
(144,108,144,36,108,144),
\tag{8.4}
\]

whereas the three central words \(122,131,221\) have

\[
(144,108,132,36,108,144).
\tag{8.5}
\]

The latter lose precisely the \(\mathbb H\otimes V\) spin profile

\[
(1,6,5,0,0)
\tag{8.6}
\]

from the long-edge image. The exceptional word \(212\) instead restores
the 144-dimensional standard core and adds the transverse cross-edge
quaternion of Section 7. Hence the reduced layer has three exact
structural classes:

\[
\boxed{
\begin{array}{c|c}
113,311&\text{standard exterior-type profile}\\
122,131,221&\text{central long-edge defect}\\
212&\text{standard core plus cross-edge }\mathbb H.
\end{array}}
\tag{8.7}
\]

In (8.6), the entries are the dimensions of the spin-\(0,1,2,3,4\)
isotypic subspaces, not the numbers of irreducible copies. They sum to
12.

Together with exterior suspension, this closes all fifteen tetrahedral
supports at \(n=7\) over \(\mathbb Q\). Quotient existence and dimension
are no longer the finite problem. Note 24 takes the next step by
comparing all transported direct-anchor completions.

\begin{center}\rule{0.5\linewidth}{0.5pt}\end{center}

\hypertarget{the-transported-anchor-transition-groupoid}{%
\subsection{9. The transported-anchor transition
groupoid}\label{the-transported-anchor-transition-groupoid}}

Write the two ordered spectator factors as \(V_a,V_b\), \(a<b\), and put

\[
T_7=\mathbb H_L\otimes\mathbb H_R\otimes V_a\otimes V_b.
\tag{9.1}
\]

If \(\phi_{\alpha,r}\) applies an operator
\(\phi:\mathbb H\otimes V\to\mathbb H\otimes V\) to the indicated
quaternion/vector pair, Note 24 proves that every transported
direct-anchor transition is generated by

\[
\boxed{
A=(\theta^{-1})_{R,a},
\qquad
B=(\theta^{-1})_{R,b},
\qquad
S=(-\theta)_{L,b}.
}
\tag{9.2}
\]

The generic history chains are

\[
\boxed{
122:S,A,
\qquad
131:B,S,A,
\qquad
221:B,S,
}
\tag{9.3}
\]

while all admissible histories collapse to one coordinate at \(113\) and
\(311\). Across all six reduced words, twenty histories yield fourteen
distinct core coordinates.

The factor supports determine the algebra:

\[
AS=SA,
\qquad
\operatorname{rank}[A,B]=64,
\qquad
\operatorname{rank}[S,B]=51.
\tag{9.4}
\]

The full internal composite has

\[
\boxed{
m_{ASB}(t)
=(t-1)^2(t+1)^3(t^2+1)(t^2-t+1).
}
\tag{9.5}
\]

Thus overlapping copies of the first cross-product shear create
\(\Phi_4\) and \(\Phi_6\) rotational factors. At the exceptional word
the full transition is instead sharply block diagonal:

\[
\boxed{
(C_1,\kappa_{212})
=
(S\oplus I_{\mathbb H})(C_0,\kappa_{212}).
}
\tag{9.6}
\]

The extra quaternion is fixed rather than converted into a core shear.
This rules out the simplest attempt to call the two-history comparison
curvature. Note 25 then constructs the required transport between
distinct spacing words and tests its only independent loop.

\begin{center}\rule{0.5\linewidth}{0.5pt}\end{center}

\hypertarget{the-spacing-word-connection-and-localized-defect}{%
\subsection{10. The spacing-word connection and localized
defect}\label{the-spacing-word-connection-and-localized-defect}}

The six positive compositions of \(5\) form the reduced word graph

\[
\Gamma_7
=
\{113\!-\!122,\ 122\!-\!131,\ 122\!-\!212,
131\!-\!221,\ 212\!-\!221,\ 221\!-\!311\}.
\tag{10.1}
\]

An elementary edge slides one internal support vertex by one gap. The
two support tetrahedra share a triangular face. Inside that face, the
outer edge on the stationary side is present as the same decoded
response variable in both local complexes.

For every adjacent pair \(\lambda,\mu\), the two selected
144-dimensional core coordinates have full rank and the same row space
on that stationary hinge. Hence the equation

\[
\boxed{
(C_\mu)_{h_{\lambda\mu}}
=g_{\lambda\mu}(C_\lambda)_{h_{\lambda\mu}}
}
\tag{10.2}
\]

defines a unique integral \(SO(3)\)-equivariant isomorphism. These six
arrows are genuine inter-object transports, not coordinate changes on
one quotient.

The graph has one independent cycle. Its exact holonomy is

\[
\boxed{
g_{212,122}g_{221,212}g_{131,221}g_{122,131}
=I_{144}.
}
\tag{10.3}
\]

Thus the complete core connection is flat. This does not make the gluing
placement-blind. On the contrary, if all three common-face blocks are
concatenated, the two trace row spaces are maximally transverse:

\[
\dim(L_{\lambda,F}\cap L_{\mu,F})=0
\tag{10.4}
\]

on every graph edge. Generic face trace spaces project isomorphically to
their common 144-dimensional hinge row space. The full exceptional trace
instead has dimension 148 and a four-dimensional projection kernel. Thus
it is an extension over that hinge, not the graph of a function on it.
These distinct face traces prevent naive facewise identification.

At the exceptional word,

\[
Y_{7,Q_{212}}\cong T_{212}\oplus K_{212},
\qquad K_{212}\cong\mathbb H.
\tag{10.5}
\]

Both graph hinges incident to \(212\) are outer edges, and the closed
square coordinate vanishes on them:

\[
(\kappa_{212})_{13}=(\kappa_{212})_{46}=0.
\tag{10.6}
\]

Using the selected splitting, define the full maps to neighboring
generic fibers by \(\rho_{212,\mu}=g_{212,\mu}\pi_T\). These maps have
kernel \(K_{212}\). Their zero action on the residual is part of this
definition: the hinge equation alone allows every extension \[
 \rho_B(t,k)=g_{212,\mu}t+B k,\qquad
 B\in\operatorname{Hom}(K_{212},T_\mu).
\] Equivariance restricts \(B\) to equivariant maps; it does not force
\(B=0\). For example, in the displayed \(T_7\) model the nonzero
injection \(q\mapsto q\otimes1\otimes\sum_a e_a\otimes e_a\) is
equivariant. Thus the intrinsic assertion is invisibility on the hinges;
forgetting is the chosen extension. With that convention the complete
finite answer is

\[
\boxed{
\text{flat transportable 144-core}
\quad+\quad
\text{\(212\)-supported quaternionic defect}.
}
\tag{10.7}
\]

The core loop is exactly flat and the defect does not enter the chosen
hinge transport. This separates chart shear, word holonomy, and
placement residual within the finite construction.

\begin{center}\rule{0.5\linewidth}{0.5pt}\end{center}

\hypertarget{three-internal-spectators-the-222-residual-and-descent-obstruction}{%
\subsection{11. Three internal spectators: the 222 residual and descent
obstruction}\label{three-internal-spectators-the-222-residual-and-descent-obstruction}}

We now use actual gap labels. Let \(Q_{222}=(1,3,5,7)\), with spectators
\(2,4,6\). Removing these spectators one at a time gives parents 122,
212, and 221, respectively. The combinatorial deletion alone does not
define a map of quotients.

\hypertarget{the-intrinsic-residual-quotient}{%
\subsubsection{11.1 The intrinsic residual
quotient}\label{the-intrinsic-residual-quotient}}

For a five-probe response define \[
 \epsilon_5(F)(z)=\sum_{a,b=1}^3F(e_a,e_a,z,e_b,e_b),
\] and put \[
 \kappa_{222}(F)=\epsilon_5(F_{15})-\epsilon_5(F_{17})
                 -\epsilon_5(F_{35})+\epsilon_5(F_{37}).
\] The blocks on 13 and 57 are zero. For \(v\in V\), the identity
\(\sum_a e_a v e_a=v\) removes the two paired probes. Each literal
cross-edge response reduces to the same response \[
 v_8v_7v_6v_5\,z\,v_4v_3v_2v_1.
\] Surjectivity of each terminal face therefore makes its two cross-edge
compositions agree. Their signed incidence contributions cancel, giving
\(\kappa_{222}\partial_{222}=0\). The three independent choices of the
open probe \(z\), each with four output components, show that the map is
onto \(\mathcal R_1\).

\textbf{Theorem 11.1 (the 222 object).} Over \(\mathbb Q\), \[
 \operatorname{rank}\partial_{222}=5388,\qquad\dim Y_{222}=444.
\] The dimensions of the labelled edge images are

\begin{longtable}[]{@{}lrrrrrr@{}}
\toprule\noalign{}
Actual edge & 13 & 15 & 17 & 35 & 37 & 57 \\
\midrule\noalign{}
\endhead
\bottomrule\noalign{}
\endlastfoot
Image dimension & 432 & 336 & 408 & 120 & 336 & 432 \\
\end{longtable}

The two outer-edge images coincide and give an intrinsic exact sequence
\[
 0\longrightarrow T_{222}\longrightarrow Y_{222}
 \xrightarrow{\bar\kappa_{222}}\mathcal R_1\longrightarrow0,
 \qquad \dim T_{222}=432.
\] In particular, \[
 T_{222}=E^Y_{13}=E^Y_{57}=\ker\bar\kappa_{222},\qquad
 D_{222}:=Y_{222}/T_{222}\cong\mathcal R_1.
\] No complementary 12-dimensional subspace inside \(Y_{222}\) is
chosen.

\emph{Proof.} The companion certificate constructs the rational matching
in natural face coordinates and right-decoded edge coordinates. Its
target dimension is 5832. A reconstructed 444-row annihilator is checked
against the matching over the integers and has full row rank. Matching
ranks 5388 over both prime fields give the opposite bound over
\(\mathbb Q\). The six edge ranks and the 432-dimensional joint
outer-edge image are verified using exact rational row identities. Since
\(\kappa_{222}\) is onto and vanishes on the outer edges, each
432-dimensional outer image equals its kernel in \(Y_{222}\). This
proves the sequence and the equalities. Full construction and all finite
matrices are reproducible from Note 26's certificate.

On the generic parent sides, Note 22's exterior edge suspension gives \[
 Y_{122}\otimes V\xrightarrow{\sim}T_{222},\qquad
 Y_{221}\otimes V\xrightarrow{\sim}T_{222}.
\] The first uses actual child edge 13 and the second uses 57 (with the
standard tensor flip for the left suspension). In each case the carried
parent block and the child block have the same kernel and rank 432. This
proves the unique descended isomorphism onto the core; it does not
define a projection from the complete child quotient onto that core.

\hypertarget{a-specified-central-slot-decoder-comparison}{%
\subsubsection{11.2 A specified central-slot decoder
comparison}\label{a-specified-central-slot-decoder-comparison}}

For the central parent, put \(D_{212}=Y_{212}/T_{212}\), where
\(T_{212}=\ker\bar\kappa_{212}\) is the common 144-dimensional
outer-edge image. Thus \(D_{212}\cong\mathbb H\). The anchored summand
\(K_{212}\) of Section 7 maps isomorphically onto this intrinsic
quotient.

Choose the \textbf{central-slot right decoder suspension}. On every
edge, insert an argument \(z\) at actual gap 4, but multiply its value
on the right: \[
 (\Sigma_4^R(F\otimes w))(\ldots,z,\ldots)
       =F(\ldots,\widehat z,\ldots)\,z\,w.
\] This is an explicitly chosen response operator. It is not identified
with inserting a vector into a literal state word, and no full chain-map
property is assumed. Invertibility of the local encoder makes it an
isomorphism \(E_{212}\otimes V\to E_{222}\).

On the cross edges, the open argument is the middle one, so the paired
sums give the exact residual square \[
 \boxed{\kappa_{222}\Sigma_4^R
 =\Theta_R(\kappa_{212}\otimes\operatorname{id}_V),
 \qquad \Theta_R(k\otimes w)(z)=kzw.}
\] This identity is checked edge by edge in the certificate as well.

\textbf{Theorem 11.2 (residual comparison and full descent
obstruction).} The chosen suspension induces an isomorphism \[
 \beta_R:D_{212}\otimes V\xrightarrow{\sim}D_{222},
 \qquad\operatorname{rank}\beta_R=12,\quad\ker\beta_R=0.
\] For a fixed unit \(w\in V\), the corresponding map on the parent
residual has rank four and zero kernel, with readout \[
 f_k(z)=kzw,\qquad k=-f_k(w).
\] However, the full suspension does not descend to
\(Y_{212}\otimes V\to Y_{222}\). Its obstruction \[
 \mathcal O=\pi_{222}\Sigma_4^R
            (\partial_{212}\otimes\operatorname{id}_V)
\] satisfies \[
 \boxed{\operatorname{im}\mathcal O=T_{222},
        \qquad\operatorname{rank}\mathcal O=432.}
\] Hence \(D_{222}\) is exactly the largest quotient of \(Y_{222}\) on
which this suspension descends from \(Y_{212}\otimes V\).

\emph{Proof.} The residual square sends old matching relations into
\(\ker\bar\kappa_{222}=T_{222}\). Passing to that quotient gives
\(\Theta_R(\bar\kappa_{212}\otimes\operatorname{id}_V)\); its kernel is
\(T_{212}\otimes V\) and its induced residual map is an isomorphism. For
fixed unit \(w\), \(w^2=-1\) gives the displayed inverse readout. The
same residual square bounds the obstruction rank by 432. The certificate
exhibits rank 432 modulo each prime, using operators whose exact
rational identities have been verified. The upper and lower bounds
agree. Every quotient admitting full descent must kill this entire
image; quotienting by it is also sufficient. This proves the final
universal property.

\hypertarget{what-has-and-has-not-been-transported}{%
\subsubsection{11.3 What has and has not been
transported}\label{what-has-and-has-not-been-transported}}

The earlier projection-defined maps forget \(K_{212}\). The operation in
this section is a different, specified comparison into a new residual
quotient; its successful readout does not reverse those
information-losing hinges. The right decoder choice is part of Theorem
11.2. A choice-independent full transport, or a local rule that turns
the representative ambiguity into a new core state, is not provided.

The finite obstruction supplies a concrete next question: which
additional local condition, if any, determines a representative or a
correction? It is not a time-update law and is not nonzero core
holonomy.

\begin{center}\rule{0.5\linewidth}{0.5pt}\end{center}

\hypertarget{results-and-scope}{%
\subsection{12. Results and scope}\label{results-and-scope}}

Stripped of the surrounding depth-filtration narrative, Notes 14--26
give the following self-contained package.

\begin{enumerate}
\def\labelenumi{\arabic{enumi}.}
\item
  \textbf{Closed tetrahedral quotient.} The minimal six-edge relation
  object is explicitly \(\mathbb H\otimes\mathbb H\), with a canonical
  \(12+4\) orthogonal channel split.
\item
  \textbf{Placement-sensitive decoration.} Isomorphic total quotients
  can be distinguished by the incidence of their six labelled edge
  images.
\item
  \textbf{Two-chart spectator atlas.} The five one-spectator placements
  require two direct seed charts, and the existence and nonexistence of
  every chart anchor are exact over \(\mathbb Q\).
\item
  \textbf{Central cross-product transition.} The exterior overlap is the
  identity, while the central overlap is the explicit
  reflection-plus-shear map \(G\).
\item
  \textbf{Channel transfer.} The central placement changes which
  orthogonal seed channel carries the surviving quaternion.
\item
  \textbf{Invariant residual.} The coordinate decoder changes across the
  overlap, but the quaternion-valued cap residual does not.
\item
  \textbf{Parity continuation.} The first admissible two-spectator cross
  square produces an exact quaternionic residual at \(n=7\).
\item
  \textbf{Exceptional core--residual decomposition.} A transported
  spectator core and the cross square give the exact split \(144+4\).
\item
  \textbf{Exterior functoriality.} Left and right exterior spectators
  tensorize every quotient and labelled edge image, and the two
  suspensions commute.
\item
  \textbf{Complete reduced two-spectator layer.} All six internal words
  have exact rational coordinates and certified edge-incidence profiles.
\item
  \textbf{Three-generator transition groupoid.} All transported-anchor
  histories are generated by pair-local \(A,B,S\), and the exceptional
  residual is fixed by its full transition.
\item
  \textbf{Stationary-edge word connection.} Every adjacent reduced word
  pair has a unique integral equivariant core transport descending from
  its shared outer hinge, and the unique independent loop is exactly
  flat.
\item
  \textbf{Localized quaternionic defect.} Full common-face traces are
  transverse, while the exceptional quaternion vanishes on both incident
  hinges and is supported only at \(212\).
\item
  \textbf{Three-spectator residual comparison.} The word 222 has an
  intrinsic 432-dimensional outer core and a 12-dimensional residual
  quotient. A specified right decoder induces an isomorphism from
  \(D_{212}\otimes V\) onto that quotient, while its full descent
  obstruction fills the core.
\end{enumerate}

Ordered placement determines which local charts exist. Their transitions
define the stationary-edge core connection, while the exceptional
residual is invisible on the two incident hinges.

This is a complete finite algebraic atlas through two spectators,
together with its exterior towers, same-word transition groupoid,
adjacent-word connection, and closed-loop test. Calling the quaternion
curvature would still be premature: the first independently defined loop
is flat, and the quaternion lies in the kernel of both selected
projection-defined maps. Section 11 settles one three-internal-spectator
object and a specified residual comparison. The remaining
three-spectator word complex and any residual-coupled holonomy remain
open.

\begin{center}\rule{0.5\linewidth}{0.5pt}\end{center}

\newpage

\hypertarget{sources-and-reproducibility}{%
\subsection{13. Sources and
reproducibility}\label{sources-and-reproducibility}}

\begin{longtable}[]{@{}
  >{\raggedright\arraybackslash}p{(\columnwidth - 2\tabcolsep) * \real{0.1000}}
  >{\raggedright\arraybackslash}p{(\columnwidth - 2\tabcolsep) * \real{0.9000}}@{}}
\toprule\noalign{}
\begin{minipage}[b]{\linewidth}\raggedright
Note
\end{minipage} & \begin{minipage}[b]{\linewidth}\raggedright
Role in this synthesis
\end{minipage} \\
\midrule\noalign{}
\endhead
\bottomrule\noalign{}
\endlastfoot
\href{../notes/14-closed-quaternionic-tetrahedral-operator.md}{14} &
closed seed operator \(\omega_5\) and the first \(12+4\) split \\
\href{../notes/15-spectator-placement-residuals.md}{15} & decorated edge
images and placement classification \\
\href{../notes/16-parity-square-quaternionic-residual.md}{16} & all-odd
parity square and exact \(n=7\) cross residual \\
\href{../notes/17-even-length-capped-five-edge-residual.md}{17} &
all-even cap complex and exact \(n=6\) residual \\
\href{../notes/18-canonical-seed-quaternion-coordinate.md}{18} &
Frobenius coordinate and orthogonal seed channels \\
\href{../notes/19-central-spectator-channel-transfer.md}{19} &
normalized \(n=6\) quotient, contraction \(\beta\), and channel
transfer \\
\href{../notes/20-n6-spectator-atlas-and-central-shear.md}{20} &
complete one-spectator atlas and central transition \\
\href{../notes/21-n7-exceptional-core-decomposition.md}{21} & exact
exceptional \(144+4\) core--residual coordinate \\
\href{../notes/22-exterior-spectator-suspension.md}{22} & all-length
exterior suspension and strict left/right flatness \\
\href{../notes/23-n7-reduced-internal-word-atlas.md}{23} & complete
rational atlas of the six reduced \(n=7\) words \\
\href{../notes/24-n7-transported-anchor-transition-groupoid.md}{24} &
exhaustive anchor histories, three transition generators, and
exceptional residual invariance \\
\href{../notes/25-n7-spacing-word-flat-core-and-quaternionic-defect.md}{25}
& stationary-edge word transport, core flatness, face transversality,
and localized quaternionic defect \\
\href{../notes/26-n8-222-residual-redetection.md}{26} & 444-dimensional
222 quotient, residual comparison, and full descent obstruction \\
\end{longtable}

The companion archive supplies the original notes and the following
certificates. Earlier certificates retain their original claim
boundaries; the release verification record identifies which were
executed for this snapshot.

\begin{Shaded}
\begin{Highlighting}[]
\NormalTok{certificates/n8\_222\_redetection\_certificate.py}
\end{Highlighting}
\end{Shaded}

The earlier exact computations are supplied by

\begin{Shaded}
\begin{Highlighting}[]
\NormalTok{certificates/n5\_quaternionic\_second\_differential\_certificate.py}
\NormalTok{certificates/n5\_central\_channel\_factorization\_certificate.py}
\NormalTok{certificates/spectator\_placement\_residual\_certificate.py}
\NormalTok{certificates/n7\_exceptional\_square\_operator\_certificate.py}
\NormalTok{certificates/n6\_capped\_five\_edge\_operator\_certificate.py}
\NormalTok{certificates/n6\_seed\_cap\_bridge\_certificate.py}
\NormalTok{certificates/n6\_spectator\_chart\_transition\_certificate.py}
\NormalTok{certificates/n7\_exceptional\_core\_decomposition\_certificate.py}
\NormalTok{certificates/exterior\_spectator\_suspension\_certificate.py}
\NormalTok{certificates/n7\_internal\_word\_atlas\_certificate.py}
\NormalTok{certificates/n7\_anchor\_transition\_groupoid\_certificate.py}
\NormalTok{certificates/n7\_spacing\_word\_transport\_certificate.py}
\end{Highlighting}
\end{Shaded}

\hypertarget{provenance-and-mathematical-dependencies}{%
\subsection{14. Provenance and mathematical
dependencies}\label{provenance-and-mathematical-dependencies}}

Section 1 gives the local definitions. Complete derivations and
computational proof components are retained in the companion notes and
certificates.

Foundation: Residual Chart Lab, \emph{Free Numbers Core v1.0.0} (2026),
DOI
\href{https://doi.org/10.5281/zenodo.21328471}{10.5281/zenodo.21328471}.

\href{https://github.com/residual-chart-lab/free-number-core}{Source
repository}: baseline \texttt{09b77f329d0c825c794443a0cd98cde89d8f1d73};
Note 26 commit \texttt{881ff66}. The claim ledger, verification records,
manifest and checksums fix the scope and reproducibility of this
MIT-licensed snapshot. No literature-wide novelty or physical
interpretation is claimed.

\end{document}
